% \textbf{Review} In the lectures, we have described the least mean square solution for linear regression as 
% \begin{equation}
%     w^* = (\tilde{X}^T \tilde{X}) \tilde{X}^T y
% \end{equation}

% where $\tilde{X}$ is the design matrix (N rows, D+1 columns) and y is the N-dimensional column vector of the true values in the training data $\mathcal{D}= \{(x_n,y_n)\}_{n=1}^N$.


% \textbf{1.1} We mentioned a practical challenge for linear regression: when $\tilde{X}^T\tilde{X}$ is not invertible. Please use a concise mathematical statement (\emph{in one sentence}) to summarize the relationship between the training data $\hat{X}$ and the dimensionality of $w$ when this bad scenario happens. Then use this statement to explain why this scenario must happen when $N < D + 1$.

We will consider the problem of fitting a linear model. Given
$d$-dimensional input data
$\textbf{x}_1,\dotsb, \textbf{x}_n \in \mathbb{R}^d$ with
real-valued labels $y_1,\dotsb, y_n \in \mathbb{R}$, the goal
is to find the coefficient vector $\textbf{w}$ that minimizes the sum
of the squared errors.  The total squared error of \textbf{w} can be
written as
%$f(\mathbf{w})=\sum_{i=1}^{n}(\mathbf{w^T x_i}-y_i)^2$,
%which is the same as the sum
$F(\mathbf{w})=\sum_{i=1}^{n}f_i(\mathbf{w})$, where
$f_i(\mathbf{w})=(\mathbf{w^T x_i}-y_i)^2$ denotes the
squared error of the $i$th data point. We will refer to $F(\w)$ as the \emph{objective function} for the problem.\\

The data in this problem will be drawn from the following linear
model.  For the training data, we select $n$ data points
$\textbf{x}_1,\dotsb, \textbf{x}_n$, each drawn independently
from a $d$-dimensional Gaussian distribution.  We then pick
the ``true'' coefficient vector $\mathbf{w}^*$ (again from a
$d$-dimensional Gaussian), and give each training point
$\mathbf{x}_i$ a label equal to $(\mathbf{w}^*)^{\top}\mathbf{x}_i$
plus some noise (which is drawn from a
1-dimensional Gaussian distribution).\\
%, and
%corresponding $y$ values $y^{(1)},\dotsb, y^{(n)}$ with $y_i$ set
%to be a linear function of $\textbf{x}_i$ with some independent
%Gaussian noise.\\ 

The following Python code can be found in your starter code and will generate the data used in this problem.

\begin{verbatim}
d = 100 # dimensions of data
n = 1000 # number of data points
X = np.random.normal(0,1, size=(n,d))
w_true = np.random.normal(0,1, size=(d,1))
y = X.dot(w_true) + np.random.normal(0,0.5,size=(n,1))
\end{verbatim}

\textbf{Part 4.1} (\blue{6pts}) Least-squares regression has the closed form solution $\textbf{w}_{LS} =
\mathbf{(X^TX)^{-1}X^Ty}$, which minimizes the squared error on the
data. 
(Here $\mathbf{X}$ is the $n \times d$ data matrix as in the code
above, with one row per data point, and $\mathbf{y}$ is the $n$-vector
of their labels.)
Solve for $\textbf{w}_{LS}$ on the training data and report the value of the objective function $F(\w_{LS})$.  For comparison, what is the 
%test 
%set error 
total squared error $F(\w)$ if you just set $\textbf{w}$ to be the all 0's vector? Using similar Python code, draw $n=1000$ test data points from the same distribution, and report the total squared error of $\w_{LS}$ on these test points. What is the gap in the training and test objective function values? Comment on the result. \\

Note: Computing the closed-form solution requires time
$O(nd^2 + d^3)$, which is slow for large $d$. Although gradient
descent methods will not yield an exact solution, they do give a
close approximation in much less time. For the purpose of this
assignment, you can use the closed form solution as a good sanity
check in the following parts.\\

\textit{Comment:} The outputs of the least-squares model are 
$F(w_{\text{LS}}^{\text{test}}) = 294.06$ and 
$F(w_{\text{LS}}^{\text{train}}) = 217$. 
The relatively small gap between the training error and the test error suggests that the model generalizes well to unseen data. 
This behavior is expected for a well-specified linear model trained with a sufficient amount of data.

\bigskip
\noindent \emph{Task}: Fill in the code for the function \colorbox{Gainsboro!60!Lavender}{closed\_form} and \colorbox{Gainsboro!60!Lavender}{square\_loss}. \\

\textbf{Part 4.2} (\blue{7pts}) In this part, you will solve the same problem via
  \emph{gradient descent} on the squared-error objective function
  $F(\textbf{w})=\sum_{i=1}^n f_i(\textbf{w})$.  Recall that the gradient of a
  sum of functions is the sum of their gradients.  Given a point
  $\textbf{w}_t$, what is the gradient of $f$ at $\textbf{w}_t$?

  Now use gradient descent to find a coefficient vector $\textbf{w}$ that 
approximately minimizes the
  least squares objective function over the training data. Run gradient descent
  three times, once with each of the step sizes
  $0.00005$, $0.0005$, and $0.0007$. You should initialize $\textbf{w}$ to be
  the all-zero vector for all three runs. Plot the objective function
  value for 20 iterations for all 3 step sizes on the same
  graph. You can also have more graphs on separate step size if the curves with all 3 step sizes are hard to read. Comment in 3-4 sentences on how the step size can affect the
  convergence of gradient descent (feel free to experiment with other
  step sizes). Also report the step size among \{0.00005, 0.0005, 0.0007\} that had the best final
  objective function value and the corresponding objective function value.\\

\textit{Comment:} From the experimental results, smaller step sizes lead to smooth but very slow decreases in the objective value across iterations. 
Moderate step sizes converge much faster and achieve a lower final squared error within the same number of iterations. 
However, when the step size becomes too large, the objective value starts to oscillate and may even increase, indicating unstable updates or divergence. 
These observations demonstrate that choosing an appropriate step size is crucial for balancing convergence speed and stability in gradient descent, as it provides the best trade-off between fast convergence and numerical stability.

\bigskip
  \noindent \emph{Task}: Fill in the code for the function \colorbox{Gainsboro!60!Lavender}{gradient\_descent}. You may want to plot the curves for a subset of the step sizes to make the curves distinguishable. \\

\textbf{Part 4.3} (\blue{7pts})
In this part you will run \emph{stochastic gradient descent} to solve the same problem.  Recall that in stochastic gradient descent, you pick one training datapoint at a time, say $(\textbf{x}_i,y_i)$, and update your current value of $\textbf{w}$ according to the gradient of $f_i(\mathbf{w})=(\mathbf{w^T x_i}-y_i)^2$. \\

Run stochastic gradient descent using step sizes
$\{0.0005, 0.005, 0.01\}$ and 1000 iterations. You should initialize $\mathbf{w}$ to be the all-zero vector for all three runs. Plot the objective
function value vs. the iteration number for all 3 step sizes on the
same graph. Comment 3-4 sentences on how the step size can affect the
convergence of stochastic gradient descent and how it compares to
gradient descent. Compare the performance of the two methods. How do
the best final objective function values compare? How many times does
each algorithm use each data point? Also report the step size that had
the best final objective function value and the
corresponding objective function value.\\

\textit{Comment:} The step size continues to play an important role in stochastic gradient descent (SGD). 
A small step size results in slow but stable progress, whereas a larger step size accelerates convergence at the cost of noticeable oscillations in the loss curve. 
When comparing the loss curves of SGD and batch gradient descent, both methods exhibit a similar overall decreasing trend. 
However, the gradient descent curve is much smoother because it uses the full dataset for each update. 
In contrast, the SGD curve is noisier due to random sampling, but it benefits from performing a much larger number of updates (1000 iterations). 
As a result, despite its higher variance, SGD achieves a significantly lower final loss, primarily because it carries out many more parameter updates.


\bigskip
\noindent \emph{Task}: Fill in the code for the function \colorbox{Gainsboro!60!Lavender}{stochastic\_gradient\_descent}. You may want to plot the curves for a subset of the step sizes to make the curves distinguishable. \\

\iffalse
\textbf{5.4} (\blue{0pts}) Include the code for all the previous parts in the submitted pdf file.
\fi

\textbf{Code submission}\\

Make sure to include your complete filled-in code in the single zip file that you upload for all your codes for the HW.